
\medskip
\noindent\SGARefTarget{sga3:viii:corollary:1:4}\textbf{Corollary 1.4.}
\textit{A diagonalizable group is reflexive, and therefore so is a locally
diagonalizable group. If \(M\) and \(N\) are ordinary commutative groups,
the natural homomorphism}
\SGARefTarget{sga3:viii:corollary:1:4:display:natural-isomorphism:01}\[
  \operatorname{Hom}_{S\text{-}\mathrm{gr}}(M_S,N_S)
  \longrightarrow
  \operatorname{Hom}_{S\text{-}\mathrm{gr}}(D(N_S),D(M_S))
\]
\textit{is bijective.}

\footnote{Editor's note: the discussion that follows has been expanded.}
\SGARefTarget{sga3:viii:corollary:1:4:discussion:01}Since \SGARefLink{sga3:viii:corollary:1:4:display:natural-isomorphism:01}{the preceding isomorphism} is compatible with base extension, it yields
an isomorphism of \(S\)-functors in groups
\SGARefTarget{sga3:viii:eq:1:4:1}
\begin{equation}
  \tag{1}
  \operatorname{\underline{Hom}}_{S\text{-}\mathrm{gr}}(M_S,N_S)
  \xrightarrow{\ \sim\ }
  \operatorname{\underline{Hom}}_{S\text{-}\mathrm{gr}}
    (D(N_S),D(M_S)).
\end{equation}
For every \(S\)-prescheme \(T\),
\[
  \operatorname{\underline{Hom}}_{S\text{-}\mathrm{gr}}(M_S,N_S)(T)
  =
  \operatorname{Hom}_{\mathrm{gr}}\bigl(M,\Gamma(N_T/T)\bigr),
\]
and by I, \SGARefLink{sga3:i:subsection:1:8}{1.8}, \(\Gamma(N_T/T)\) is the abelian group of locally constant
maps \(T\to N\). On the other hand, let
\(\operatorname{Hom}_{\mathrm{gr}}(M,N)_S\) be the constant \(S\)-group
associated with the ordinary abelian group
\(\operatorname{Hom}_{\mathrm{gr}}(M,N)\). There is an evident
homomorphism of \(S\)-functors in commutative groups
\SGARefTarget{sga3:viii:eq:1:4:2}
\begin{equation}
  \tag{2}
  \operatorname{Hom}_{\mathrm{gr}}(M,N)_S
  \xrightarrow{\ \theta\ }
  \operatorname{\underline{Hom}}_{S\text{-}\mathrm{gr}}(M_S,N_S),
\end{equation}
which is always a monomorphism and is an isomorphism if \(M\) is of finite
type.\footnote{Editor's note: write \(F(M)\) and \(G(M)\) for the left- and
right-hand sides of \SGARefLink{sga3:viii:eq:1:4:2}{(2)}, respectively. If \(M=M_1\oplus M_2\), there are
canonical isomorphisms
\(F(M)=F(M_1)\oplus F(M_2)\) and
\(G(M)=G(M_1)\oplus G(M_2)\). It is therefore enough to check that
\(\theta\) is an isomorphism for \(M=\mathbf Z/r\mathbf Z\), \(r\geq0\).
In that case \(F(M)=({}_rN)_S\), where \({}_rN\) is the kernel of
\(r\operatorname{id}_N\), and for every \(T\to S\) the homomorphism
\[
  F(M)(T)=\Gamma({}_rN_T/T)
  \longrightarrow
  G(M)(T)={}_r\Gamma(N_T/T)
\]
is bijective.}

Part \SGARefLink{sga3:viii:corollary:1:5:item:a}{(a)} of \SGARefLink{sga3:viii:corollary:1:5}{the following corollary} follows from \SGARefLink{sga3:viii:corollary:1:4:discussion:01}{the preceding discussion};
part \SGARefLink{sga3:viii:corollary:1:5:item:b}{(b)} follows from it by the descent results recalled in \SGARefLink{sga3:viii:paragraph:1-7}{1.7}.
\footnote{Editor's note: the original stated after \SGARefLink{sga3:viii:corollary:1:5}{1.5} that, more generally,
if \(G\) and \(H\) are locally diagonalizable and \(H\) is of finite type,
then \(\operatorname{\underline{Hom}}_{S\text{-}\mathrm{gr}}(G,H)\) is
representable. \SGARefLink{sga3:viii:corollary:1:5:item:b}{That assertion} has been included in \SGARefLink{sga3:viii:corollary:1:5}{1.5} and its proof made
explicit in \SGARefLink{sga3:viii:paragraph:1-7}{1.7}.}

\medskip
\noindent\SGARefTarget{sga3:viii:corollary:1:5}\textbf{Corollary 1.5.}
\begin{enumerate}
  \item[(a)]\SGARefTarget{sga3:viii:corollary:1:5:item:a}
  \textit{Let \(M\) and \(N\) be ordinary commutative groups, with \(M\) of
  finite type. Then there is an isomorphism}
  \[
    \operatorname{Hom}_{\mathrm{gr}}(M,N)_S
    \xrightarrow{\ \sim\ }
    \operatorname{\underline{Hom}}_{S\text{-}\mathrm{gr}}
      (D(N_S),D(M_S));
  \]
  \textit{consequently,
  \(\operatorname{\underline{Hom}}_{S\text{-}\mathrm{gr}}
  (D(N_S),D(M_S))\) is representable.}

  \item[(b)]\SGARefTarget{sga3:viii:corollary:1:5:item:b}
  \textit{More generally, if \(G\) and \(H\) are locally diagonalizable and
  \(H\) is of finite type, then
  \(\operatorname{\underline{Hom}}_{S\text{-}\mathrm{gr}}(G,H)\) is
  representable.}
\end{enumerate}

\footnote{Editor's note: a statement 1.5.1 treating the functor
\(\operatorname{\underline{Isom}}_{S\text{-}\mathrm{gr}}(G,H)\), considered
in X, \SGARefLink{sga3:x:corollary:5-10}{5.10} and 5.11, should be added.}
Corollary~\SGARefLink{sga3:viii:corollary:1:5}{1.5} gives:

\medskip
\noindent\SGARefTarget{sga3:viii:corollary:1:6}\textbf{Corollary 1.6.}
\textit{Under the hypotheses of \SGARefLink{sga3:viii:corollary:1:5}{1.5}, if \(S\) is connected, then}
\[
  \operatorname{Hom}_{S\text{-}\mathrm{gr}}(D_S(N),D_S(M))
  \xrightarrow{\ \sim\ }
  \operatorname{Hom}_{\mathrm{gr}}(M,N)
\]
\textit{and}
\[
  \operatorname{Isom}_{S\text{-}\mathrm{gr}}(D_S(N),D_S(M))
  \xrightarrow{\ \sim\ }
  \operatorname{Isom}_{\mathrm{gr}}(M,N).
\]

\medskip
\noindent\SGAInlineRefTarget{sga3:viii:paragraph:1-7}\SGARefTarget{sga3:viii:subsection:1:7}\textbf{1.7. Descent of
representability.}
\footnote{Editor's note: \SGARefLink{sga3:viii:paragraph:1-7}{this paragraph} has been added to make explicit the
``local on \(S\)'' character of the representability of certain sheaves on
\(S\), which is used repeatedly below and was used implicitly in the
original.}
We recall here several descent results that will be used frequently below.

\medskip
\noindent\SGARefTarget{sga3:viii:scholium:1:7:1}\textbf{Scholium 1.7.1.}
\textit{Let \(S\) be a prescheme and let \(X,Y,T\) be \(S\)-preschemes. If
\((T_i)\) is an open covering of \(T\), and if
\(T_{ij}=T_i\cap T_j=T_i\times_TT_j\), then, because giving a morphism of
\(T\)-preschemes \(X_T\to Y_T\) is local on \(T\), there is an exact
sequence of sets}
\SGARefTarget{sga3:viii:eq:1:7:1}
\begin{equation}
  \tag{1}
  \operatorname{Hom}_T(X_T,Y_T)
  \longrightarrow
  \prod_i\operatorname{Hom}_{T_i}(X_{T_i},Y_{T_i})
  \rightrightarrows
  \prod_{i,j}\operatorname{Hom}_{T_{ij}}(X_{T_{ij}},Y_{T_{ij}}).
\end{equation}
\textit{Thus \(\operatorname{\underline{Hom}}_S(X,Y)\) is a local
\(S\)-functor, that is, a sheaf on \((\mathbf{Sch})_{/S}\) for the Zariski
topology.}

More generally, by IV, \SGARefLink{sga3:iv:corollary:4:5:13}{4.5.13},
\(\operatorname{\underline{Hom}}_S(X,Y)\) is a sheaf on
\((\mathbf{Sch})_{/S}\) for every topology coarser than the canonical
topology, for example the fpqc topology. If \(G\) and \(H\) are group
preschemes over \(S\), it follows that the subfunctor
\(\operatorname{\underline{Hom}}_{S\text{-}\mathrm{gr}}(G,H)\) is an fpqc
sheaf, hence a fortiori a local functor.

\medskip
\noindent\SGARefTarget{sga3:viii:lemma:1:7:2}\textbf{Lemma 1.7.2.}
\footnote{Editor's note: see also Remark XI, \SGARefLink{sga3:xi:remark:3:4}{3.4}.}
\textit{Let \(F\) be a local \(S\)-functor.}
\begin{enumerate}
  \item[(i)]\SGARefTarget{sga3:viii:lemma:1:7:2:item:i}
  \textit{Suppose that there is an open covering \((S_i)\) of \(S\) such
  that the restriction \(F_i=F\times_SS_i\) of \(F\) to each \(S_i\) is
  representable by an \(S_i\)-prescheme \(X_i\). Then \(F\) is
  representable by an \(S\)-prescheme \(X\).}

  \item[(ii)]\SGARefTarget{sga3:viii:lemma:1:7:2:item:ii}
  \textit{Suppose that \(F\) is an fpqc sheaf and that there is a faithfully
  flat quasi-compact morphism \(S'\to S\) such that
  \(F'=F\times_SS'\) is representable by an \(S'\)-prescheme \(X'\). Then
  \(X'\) carries a descent datum relative to \(S'\to S\) (cf. IV, \SGARefLink{sga3:iv:definition:2:1}{2.1}).
  If this descent datum is effective---as it is when \(X'\) is affine over
  \(S'\)---then \(F\) is representable by an \(S\)-prescheme \(X\).}
\end{enumerate}

\emph{Proof.}
\SGARefLink{sga3:viii:lemma:1:7:2:item:i}{(i)} The hypothesis implies that \(X_i\times_SS_j\) and
\(X_j\times_SS_i\) both represent the restriction of \(F\) to
\(S_{ij}=S_i\times_SS_j\). The Yoneda lemma therefore gives a unique
isomorphism of \(S_{ij}\)-preschemes
\[
  c_{ji}:X_i\times_SS_j\xrightarrow{\ \sim\ }X_j\times_SS_i.
\]
Over \(S_{ijk}=S_i\times_SS_j\times_SS_k\), we obtain the commutative
diagram of isomorphisms
\SGARefTarget{sga3:viii:diagram:1:7:2:i}
\[
\begin{tikzcd}[column sep=small,row sep=large]
X_i\times_SS_j\times_SS_k
  \arrow[r,"c_{ji}\times\operatorname{id}_{S_k}"]
  \arrow[d,equal]
&
X_j\times_SS_i\times_SS_k
  \arrow[r,equal]
&
X_j\times_SS_k\times_SS_i
  \arrow[d,"c_{kj}\times\operatorname{id}_{S_i}"]
\\
X_i\times_SS_k\times_SS_j
  \arrow[r,"c_{ki}\times\operatorname{id}_{S_j}"]
&
X_k\times_SS_i\times_SS_j
  \arrow[r,equal]
&
X_k\times_SS_j\times_SS_i .
\end{tikzcd}
\]
Indeed, all of its objects represent the restriction of \(F\) to
\(S_{ijk}\). Thus the \(c_{ji}\) satisfy the usual cocycle relation
\(c_{kj}\circ c_{ji}=c_{ki}\).

It follows that the \(X_i\) glue to an \(S\)-prescheme \(X\) with
\(X\times_SS_i=X_i\) for every \(i\). For every \(Y\) over \(S_i\),
\SGARefTarget{sga3:viii:eq:1:7:2:star}
\begin{equation}
  \tag{\(\ast\)}
  F(Y)=F_i(Y)
  =\operatorname{Hom}_{S_i}(Y,X\times_SS_i)
  =\operatorname{Hom}_S(Y,X)=h_X(Y).
\end{equation}
For an arbitrary \(Y\to S\), the \(Y_i=Y\times_SS_i\) form an open covering
of \(Y\); put \(Y_{ij}=Y_i\times_YY_j=Y\times_SS_{ij}\). Since \(F\) is
local by hypothesis and \(h_X\) is local because the Zariski topology is
coarser than the canonical topology, both \(F(Y)\) and \(h_X(Y)\), in view
of \SGARefLink{sga3:viii:eq:1:7:2:star}{\((\ast)\)}, identify with the
kernels of the respective double arrows
\[
  \prod_iF(Y_i)\rightrightarrows\prod_{i,j}F(Y_{ij}),
  \qquad
  \prod_ih_X(Y_i)\rightrightarrows\prod_{i,j}h_X(Y_{ij}).
\]
This proves \SGARefLink{sga3:viii:lemma:1:7:2:item:i}{(i)}.

\SGARefLink{sga3:viii:lemma:1:7:2:item:ii}{(ii)} Put \(S''=S'\times_SS'\), and let \(S''_1\) and \(S''_2\) denote
\(S''\) regarded as an \(S'\)-prescheme through the first and second
projections. The functors \(F''_1=F'\times_{S'}S''_1\) and
\(F''_2=F'\times_{S'}S''_2\) are represented respectively by
\(X''_1=X'\times_{S'}S''_1\) and \(X''_2=X'\times_{S'}S''_2\). Since
\[
  F''_1=F\times_SS''=F''_2,
\]
there is a unique \(S''\)-isomorphism
\[
  c:X''_1\xrightarrow{\ \sim\ }X''_2.
\]
Let \(q_i\), respectively \(p_{ji}\), be the projection from
\(S'''=S'\times_SS'\times_SS'\) onto its \(i\)-th factor, respectively
onto its \(i\)-th and \(j\)-th factors. Let
\(X'''_i=X'\times_{S'}S'''_i\), where \(S'''_i\) means \(S'''\) viewed as
an \(S'\)-prescheme through \(q_i\), and let
\(p_{ji}^{*}(c):X'''_i\xrightarrow{\sim}X'''_j\) be obtained from \(c\)
by base change. We obtain the diagram
\SGARefTarget{sga3:viii:diagram:1:7:2:ii}
\[
\begin{tikzcd}[column sep=large,row sep=large]
X'''_1
  \arrow[r,"p_{21}^{*}(c)"]
  \arrow[dr,"p_{31}^{*}(c)"']
&
X'''_2
  \arrow[d,"p_{32}^{*}(c)"]
\\
&X'''_3 .
\end{tikzcd}
\]
All three objects represent the restriction of \(F\) to \(S'''\), so the
diagram commutes:
\[
  p_{32}^{*}(c)\circ p_{21}^{*}(c)=p_{31}^{*}(c).
\]
Thus \(c\) is a descent datum on \(X'\) relative to \(S'\to S\)
(cf. IV, \SGARefLink{sga3:iv:definition:2:1}{2.1}).

Suppose that this descent datum is effective, so there is an \(S\)-prescheme
\(X\) such that \(X'\simeq X\times_SS'\). This is the case when \(X'\) is
affine over \(S'\), by SGA~1, VIII, \SGARefLink{sga3:viii:proposition:2:1}{2.1}.\footnote{Editor's note: for another
effectivity criterion, see X, \SGARefLink{sga3:x:lemma:5-4}{5.4}--\SGARefLink{sga3:viii:corollary:5:6}{5.6}.} For every \(Y\to S'\), one then has
\SGARefTarget{sga3:viii:eq:1:7:2:double-star}
\begin{equation}
  \tag{\(\ast\ast\)}
  F(Y)=F'(Y)
  =\operatorname{Hom}_{S'}(Y,X\times_SS')
  =\operatorname{Hom}_S(Y,X)=h_X(Y).
\end{equation}
For arbitrary \(Y\to S\), put \(Y'=Y\times_SS'\) and
\(Y''=Y'\times_YY'\simeq Y\times_SS''\). Like \(S'\to S\), the morphism
\(Y'\to Y\) is faithfully flat and quasi-compact, hence an
\(\mathcal M\)-effective epimorphism, where \(\mathcal M\) is the family of
faithfully flat quasi-compact morphisms. Thus the equivalence relation
\[
  Y'\times_YY'\rightrightarrows Y'
\]
is \(\mathcal M\)-effective and has quotient \(Y\). Since \(F\) is an fpqc
sheaf by hypothesis, while \(h_X\) is one because the fpqc topology is
coarser than the canonical topology, both \(F(Y)\) and \(h_X(Y)\), in view
of \SGARefLink{sga3:viii:eq:1:7:2:double-star}{\((\ast\ast)\)}, identify
with the respective kernels of
\[
  F(Y')\rightrightarrows F(Y'\times_YY'),
  \qquad
  h_X(Y')\rightrightarrows h_X(Y'\times_YY').
\]
This proves \SGARefLink{sga3:viii:lemma:1:7:2:item:ii}{(ii)}. \qed

\medskip
\noindent\SGARefTarget{sga3:viii:corollary:1:7:3}\textbf{Corollary 1.7.3.}
\textit{Let \(F\) be an fpqc sheaf on \((\mathbf{Sch})_{/S}\). Suppose
there is an open covering \((S_i)\) of \(S\), and for each \(i\) a
faithfully flat quasi-compact morphism \(S'_i\to S_i\), such that
\(F'_i=F\times_SS'_i\) is representable by an \(S'_i\)-prescheme \(X'_i\)
affine over \(S'_i\). Then \(F\) is representable by an \(S\)-prescheme
\(X\) affine over \(S\), with \(X\times_SS'_i=X'_i\) for every \(i\). If,
moreover, every \(X'_i\to S'_i\) is a closed immersion (respectively, a
finite \'etale morphism), then so is \(X\to S\).}

The first assertion follows from
\SGARefLink{sga3:viii:lemma:1:7:2}{1.7.2}. For the second, it is enough to
check that every morphism \(X\times_SS_i\to S_i\) is a closed immersion
(respectively, finite and \'etale); this follows from EGA~IV\(_2\), 2.7.1
(respectively, together with IV\(_4\), 17.7.3).

\medskip
\noindent\SGARefTarget{sga3:viii:remark:1:7:4}\textbf{Remark 1.7.4.}
Assertion \SGARefLink{sga3:viii:corollary:1:5:item:b}{1.5(b)} follows,
as announced, from \SGARefLink{sga3:viii:scholium:1:7:1}{1.7.1} and
\SGARefLink{sga3:viii:lemma:1:7:2:item:i}{1.7.2(i)}.
